\section{Conclusion}
\label{sec:conclusion}

We introduced Holographic Polar Arithmetic (HPA) as a self-contained toolkit.

\begin{itemize}
    \item From the $\mathbb{F}_1$-inspired multiplicative skeleton we defined a radial character $\rho$ and a multiplicative phase $\theta_\times$, yielding an embedding $\mathcal{Z}(n)=\rho(n)\e^{\iu\theta_\times(n)}$ that linearizes multiplication.
    \item We proved a no-go theorem: no nontrivial holographic embedding into $\CC^*$ can preserve linear addition together with multiplication. Additive readout therefore appears as a lossy projection from vector synthesis in $\CC$ back to a discrete coordinate lattice, with an intrinsic quantum gap.
    \item We introduced a genuine unitary scan operator $\Theta$ (irrational-rotation Koopman operator), defined time as iteration count, and constructed an orbit-wise cyclic calculus (difference, trace, finite part).
    \item Via binary window readout we obtained Sturmian/Fibonacci structures and, in the golden branch, a canonical coordinate system given by Zeckendorf decompositions (Ostrowski degeneration).
    \item We packaged three reusable modules: singularity compactification, finite-part extraction, and Zeckendorf-coded lattice path sums.
\end{itemize}

Under the minimal-complexity criteria, the golden branch is singled out not by numerology but by two concrete extremal features: Hurwitz anti-resonance and binary-minimal canonical encoding (Proposition~\ref{prop:golden_binary_minimal}). Future work can focus on robustness bounds, explicit error control for finite parts, and interfaces to multiplicative/spectral statistics.
